% Seccion 4 de la rubrica no la exige, pero el analisis exploratorio es lo
% que justifica el ACP y la eleccion del umbral. Aloja como \subsection a
% eda_cualitativas.tex. Las tablas exhaustivas de estadisticos viven en el
% diccionario (anexo); aqui solo lo que se interpreta.

\section{Análisis exploratorio}
\label{sec:eda}

\subsection{Variables cuantitativas}

\subsubsection{Cuartiles y límites de outliers}

Los cuartiles y límites de $1.5\times$IQR de las variables horarias se
reportan en sus unidades originales (\cref{tab:cuartiles-iqr}) y son los
mismos que usó la limpieza como criterio de outliers. \texttt{RAINF} y
\texttt{WDR} quedan fuera: \texttt{WDR} es circular y \texttt{RAINF} está
inflada de ceros, así que su IQR no es interpretable como límite
(\cref{sec:dic:meteorologicos}). Dos lecturas importan para lo que sigue: la
mediana de la componente zonal del viento es negativa, es decir, el viento
predominante entra del oriente (\cref{sec:contexto-estaciones}); y el límite
superior de \texttt{O3\_8h} (66.22 ppb) queda por encima del umbral de
51 ppb, de modo que la excedencia no es un outlier sino parte del rango
habitual de la variable.

\begin{table}[t]
  \centering
  \caption{Cuartiles y límites de outlier (criterio $1.5\times$IQR) de las
    variables cuantitativas horarias, en sus unidades originales.}
  \label{tab:cuartiles-iqr}
  \small
  \begin{tabular}{lrrrr}
    \toprule
    Variable & Q1 & Q3 & Lím. inf. & Lím. sup. \\
    \midrule
    \texttt{O3}       &    13.00 &    37.00 &  $-23.00$ &    73.00 \\
    \texttt{CO}       &     0.71 &     1.83 &   $-0.96$ &     3.50 \\
    \texttt{NO2}      &     6.90 &    19.70 &  $-12.30$ &    38.90 \\
    \texttt{NOX}      &    10.95 &    30.70 &  $-18.68$ &    60.33 \\
    \texttt{PM10}     &    35.00 &    74.00 &  $-23.50$ &   132.50 \\
    \texttt{PM2.5}    &    10.00 &    26.74 &  $-15.11$ &    51.85 \\
    \texttt{TOUT}     &    18.49 &    28.39 &      3.64 &    43.24 \\
    \texttt{PRS}      &   709.70 &   721.70 &    691.70 &   739.70 \\
    \texttt{WSR}      &     4.70 &    11.60 &   $-5.65$ &    21.95 \\
    \texttt{O3\_8h}   &    15.44 &    35.75 &  $-15.03$ &    66.22 \\
    \texttt{viento\_u}&  $-8.56$ &     0.21 &  $-21.73$ &    13.37 \\
    \texttt{viento\_v}&  $-3.38$ &     3.98 &  $-14.41$ &    15.01 \\
    \bottomrule
  \end{tabular}
\end{table}

La \cref{fig:boxplot-medidas} muestra la distribución de cada variable en
unidades originales. Todas las concentraciones de contaminantes presentan
cola derecha larga, con outliers que multiplican varias veces el tercer
cuartil; es la asimetría que motivó Box-Cox en \cref{sec:transformaciones}.
Las meteorológicas son más simétricas, salvo la velocidad del viento.

\begin{figure}[t]
    \centering
    \includegraphics[width = \linewidth]{boxplots_medidas.pdf}
    \caption{Diagramas de caja de cada variable cuantitativa horaria, en
      unidades originales.}
    \label{fig:boxplot-medidas}
\end{figure}

\subsubsection{Correlación}

A resolución horaria y sobre la escala transformada, los pares más fuertes
son \texttt{NO2}--\texttt{NOX} ($r=0.89$), multicolinealidad esperada entre
precursores del mismo ciclo de reacción; \texttt{O3}--\texttt{O3\_8h}
($r=0.69$), por construcción; \texttt{PM10}--\texttt{PM2.5} ($r=0.62$),
material particulado de origen común; \texttt{O3}--\texttt{TOUT} ($r=0.52$)
y \texttt{O3}--\texttt{WSR} ($r=0.43$), consistentes con el mecanismo
fotoquímico; y \texttt{O3}--\texttt{NOX} ($r=-0.43$), que refleja el consumo
de \ce{O3} por titulación con \ce{NO} cerca de fuentes de tráfico.

El ACP corre sobre la tabla diaria y necesita su propia matriz: agregar a
día cambia el tamaño de muestra y, típicamente, sube la correlación al
cancelar ruido de medición. Se calcula sobre las 16 variables continuas del
conjunto diario transformado, en el subconjunto con las 16 simultáneamente
completas ($n=15{,}856$ de 26,691 días-estación, 59.4\,\%), la misma regla
de descarte por lista que usa el modelado y que garantiza una matriz
semidefinida positiva (\cref{fig:mapa-calor-diario}).

\begin{figure}[t]
    \centering
    \includegraphics[width = \linewidth]{mapa_calor_correlacion_diario.pdf}
    \caption{Mapa de calor de correlación de Pearson entre las 16 variables
      continuas del conjunto diario transformado, $n=15{,}856$ días-estación
      completos.}
    \label{fig:mapa-calor-diario}
\end{figure}

De los pares horarios, tres tienen análogo directo en el conjunto diario:
\texttt{NO2}--\texttt{NOX} no cambia (0.89), \texttt{PM10}--\texttt{PM2.5}
sube a 0.70 y \texttt{viento\_u}--\texttt{WSR} sube en magnitud a $-0.48$.
El orden sí cambia: dominan \texttt{NO}--\texttt{NOX} ($r=0.95$) y
\texttt{RH\_media}--\texttt{RH\_min\_diurna} ($r=0.89$), inexistentes en el
mapa horario porque \texttt{NO} y \texttt{RH} no formaban parte de ese
conjunto. El primero es la redundancia definicional \texttt{NOX} $\approx$
\texttt{NO} + \texttt{NO2}; el segundo, y también
\texttt{WSR\_media}--\texttt{WSR\_min} ($r=0.71$), es la misma variable
resumida dos veces sobre ventanas distintas del día. Estas redundancias de
construcción son las que el análisis factorial (\cref{sec:analisis-factorial})
termina por expulsar.

Los diagnósticos previos al ACP confirman la multicolinealidad
(\cref{tab:diagnostico-pca}): el determinante es cercano a 0 y el número de
condición es alto. La prueba de esfericidad de Bartlett rechaza que la
matriz sea la identidad ($\chi^2(120)=169{,}140.7$, $p<0.001$), y el KMO
total libra el umbral convencional de 0.6; a nivel variable,
\texttt{RH\_media} (0.587), \texttt{PRS\_media} (0.589) y \texttt{TOUT\_max}
(0.596) quedan justo debajo sin arrastrar el total. Los cuatro diagnósticos
justifican aplicar ACP sobre los predictores diarios.

\begin{table}[t]
  \centering
  \caption{Diagnósticos previos al ACP sobre la matriz de correlación
    diaria ($n=15{,}856$, 16 variables).}
  \label{tab:diagnostico-pca}
  \small
  \begin{tabular}{lr}
    \toprule
    Diagnóstico          & Valor \\
    \midrule
    Determinante          & $2.32\times10^{-5}$ \\
    Número de condición   & 322.8 \\
    KMO total             & 0.660 \\
    \bottomrule
  \end{tabular}
\end{table}
